\documentclass[10pt,a4paper]{amsart}
\usepackage[margin=19mm]{geometry}
\usepackage{amsmath,amssymb,mathtools}
\usepackage[scr=rsfs,cal=euler]{mathalfa}
\usepackage{xfrac}
\usepackage[unicode,hidelinks,bookmarks=false]{hyperref}
\newcommand{\ie}{i.e.\ }
\newcommand{\cf}{cf.\ }
\newcommand{\resp}{respectively\ }
\newcommand{\Subsection}{\subsection}
\providecommand{\nobreakdash}{\nobreak\hbox{-}}
\setlength{\emergencystretch}{2em}
\multlinegap0pt
\usepackage{amssymb}
\usepackage{mathtools}
\usepackage[shortlabels]{enumitem}
\usepackage{xcolor}
\newtheorem{lemma}{Lemma}
\newcommand{\Kfun}{\mathcal K}
\title{The Dirichlet spectrum with respect to $L_1$ norm is $\left[\frac12,1\right]$}
\author{Nikita Shulga}
\date{Source: arXiv:2606.08865v1 (CC BY 4.0)}
\begin{document}
\maketitle

\noindent\emph{Source and licence.} The Dirichlet spectrum with respect to $L_1$ norm is $\left[\frac12,1\right]$. Version arXiv:2606.08865v1 (CC BY 4.0), distributed by arXiv under the Creative Commons Attribution 4.0 International licence, http://creativecommons.org/licenses/by/4.0/, as recorded in the arXiv licence metadata for this exact version and shown on the abstract page https://arxiv.org/abs/2606.08865v1. Full licence notice: \texttt{licenses/arxiv-2606.08865-CC-BY-4.0.txt}.\par\smallskip

\noindent\textit{Editorial scope.} Counted unit: the article's lemma lem:arithmetic-bound (source0.tex lines 591--622). The article's complete original proof of it is delivered with it (source0.tex lines 1788--1959, one \texttt{proof} environment of 172 lines, which the article places away from the statement). No mathematics is written, repaired or re-derived: the statement and the proof are the article's own lines, in the article's own notation, and no retained line is substituted or re-flowed.  Retained uncounted as the source-local context the proof needs: lines 428--431.  Nothing was omitted from any retained span, so the package carries no omission record.  No retained line contains a citation, so the wrapper carries no bibliography and no bibliography entry.  No retained line contains a figure environment, an image-inclusion command or drawing code, so no image is delivered and the package declares no asset dependency.  Editorial formatting only, in addition: an AMS \texttt{amsart} preamble that restates the article's own declarations and macros used by the retained lines, namely the environments \texttt{lemma} on separate counters, together with the standard packages those lines need.  The retained environments print in this document as Lemma 1 in order of appearance, whereas the article prints the same items with the numbering of its own full document.  The complete native source of the article as distributed in the arXiv e-print for this exact version is bundled unchanged as \texttt{sources/202365/source0.tex}.
\begin{lemma}\label{lem:endpoint-principle}
Let $I,J\subset[0,\infty)$ be compact intervals.  The maximum of
$\Kfun$ on $I\times J$ is attained at one of the four corners.
\end{lemma}

\begin{lemma}\label{lem:arithmetic-bound}
Let $b_1\ge2$ and $c_1\geq b_1(b_1-1)$ be integers.  Define
$$
        L=\frac{1}{b_1}+\frac1{c_1+1}.
$$
If $(b_1,c_1)=(2,2)$, put $U=1$.  Otherwise put
$$
        U=\frac{b_1^2}{c_1(c_1+1)-b_1-b_1^2}.
$$
Also put
$$
        V_{b_1}=\frac{c_1+1}{b_1c_1+b_1-1},
        \qquad
        V_{c_1}=\frac{b_1+1}{b_1c_1+c_1-1}.
$$
Then
\begin{equation}\label{eq:arithmetic-main}
        \Kfun(u,v)\le L
\end{equation}
whenever
\begin{equation}\label{eq:arithmetic-first-range}
        0\le u\le U,
        \qquad
        \frac{1}{b_1}\le v\le V_{b_1},
\end{equation}
and also whenever
\begin{equation}\label{eq:arithmetic-second-range}
        0\le v\le U,
        \qquad
        \frac{1}{c_1}\le u\le V_{c_1}.
\end{equation}
\end{lemma}

\begin{proof}[Proof of Lemma \ref{lem:arithmetic-bound}]
We prove the first assertion first.  By Lemma \ref{lem:endpoint-principle}, it
is enough to check the four corners of the rectangle
$[0,U]\times[1/b_1,V_{b_1}]$.  At the two corners with $u=0$, we have
$$
        \Kfun(0,v)=v\le V_{b_1}<L,
$$
where the last inequality follows from
$$
        L-V_{b_1}
        =\frac1{c_1+1}-\frac1{b_1(b_1c_1+b_1-1)}>0.
$$
Thus, it remains to prove
\begin{equation}\label{eq:corner-U-b}
        \Kfun\left(U,\frac{1}{b_1}\right)\le L,
\end{equation}
and
\begin{equation}\label{eq:corner-U-Vb}
        \Kfun(U,V_{b_1})\le L.
\end{equation}

The exceptional case $(b_1,c_1)=(2,2)$ is immediate:
$$
        L=\frac56,
        \qquad
        U=1,
        \qquad
        V_{b_1}=\frac35,
$$
and a direct substitution gives
$$
        \Kfun\left(1,\frac12\right)=\frac{11}{14}<\frac56,
        \qquad
        \Kfun\left(1,\frac35\right)=\frac{74}{95}<\frac56.
$$
Since $\Kfun$ is symmetric and $V_{c_1}=V_{b_1}=3/5$ in this exceptional case,
the second range \eqref{eq:arithmetic-second-range} follows from the same corner checks.  Assume now that
$(b_1,c_1)\ne(2,2)$, so that the displayed formula for $U$ is valid.

For $b_1=2$ and $c_1\geq3$, after substituting $U=4/(c_1(c_1+1)-6)$, inequalities
\eqref{eq:corner-U-b} and \eqref{eq:corner-U-Vb} reduce respectively to
$$
        \frac{3c_1^4+6c_1^3-17c_1^2-52c_1-20}{(c_1-2)(c_1+1)(c_1+3)(3c_1^2+3c_1-2)}>0
$$
and
$$
        \frac{(3c_1^2+5c_1+8)(3c_1^4+4c_1^3-17c_1^2-38c_1-16)}{2c_1^2(c_1+1)(c_1+3)(c_1-2)(2c_1+1)(3c_1+5)}>0,
$$
which holds for all $c_1\ge3$.

For $b_1=3$ and $c_1\geq6$, the substitution $U=9/(c_1(c_1+1)-12)$ reduces \eqref{eq:corner-U-b} and \eqref{eq:corner-U-Vb} respectively to
$$
        \frac{4c_1^4-4c_1^3-71c_1^2-162c_1-63}{(c_1-3)(c_1+1)(c_1+4)(2c_1-1)(2c_1+3)}>0
$$
and
$$
        \frac{2(4c_1^2+7c_1+9)(4c_1^4-7c_1^3-68c_1^2-120c_1-54)}{3c_1^2(c_1-3)(c_1+1)(c_1+4)(3c_1+2)(4c_1+7)}>0,
$$
which holds for all $c_1\geq6$.

It remains to treat $b_1\ge4$.  Put
$$
        c_1=b_1(b_1-1)+T,
        \qquad T\ge0.
$$
Then \eqref{eq:corner-U-b} and \eqref{eq:corner-U-Vb} reduce respectively to
$$
        \frac{P_1(T)}{D_1(T)}>0,
        \qquad
        \frac{P_2(T)}{D_2(T)}>0,
$$
where
\begin{align*}
P_1(T)={}&b_1(b_1+1)T^4
        +b_1(b_1+1)(3b_1^2-2b_1+2)T^3\\
        &+b_1(3b_1^5-3b_1+1)T^2
        +b_1^2(b_1^6+b_1^5-3b_1^4-b_1^3-3b_1^2-b_1-2)T\\
        &+b_1^6(b_1^3-3b_1^2+2b_1-3),\\[3pt]
D_1(T)={}&b_1(T+b_1^2+1)(T+b_1^2-2b_1)(T+b_1^2-b_1+1)\\
        &\times\bigl((b_1+1)T^2+(2b_1^3-b_1+1)T+b_1^5-b_1^4+b_1^2-2b_1\bigr),
\end{align*}
and
\begin{align*}
P_2(T)={}&(b_1-1)(b_1+1)^2T^6
        +(b_1-1)(b_1+1)^2(5b_1^2-4b_1+3)T^5\\
        &+(10b_1^7-5b_1^6-10b_1^5+14b_1^4-8b_1^3-3b_1^2+9b_1-3)T^4\\
        &+(10b_1^9-10b_1^8-10b_1^7+26b_1^6-30b_1^5+16b_1^4+11b_1^3\\
        &\qquad -13b_1^2+11b_1-1)T^3\\
        &+b_1(b_1-1)(5b_1^9-15b_1^7+24b_1^6-33b_1^5+18b_1^4-9b_1^2\\
        &\qquad +10b_1-4)T^2\\
        &+b_1^2(b_1^2+1)(b_1^9+b_1^8-17b_1^7+39b_1^6-49b_1^5+37b_1^4\\
        &\qquad -12b_1^3-10b_1^2+9b_1-5)T\\
        &+b_1^3(b_1^2+1)^2(b_1^3-2b_1^2+b_1+1)(b_1^4-4b_1^3+5b_1^2-5b_1+2),\\[3pt]
D_2(T)={}&b_1(T+b_1^2+1)(T+b_1^2-2b_1)(T+b_1^2-b_1)^2(T+b_1^2-b_1+1)\\
        &\times\bigl((b_1+1)T+b_1^3+b_1+1\bigr)(b_1T+b_1^3-b_1^2+b_1-1).
\end{align*}
All factors in $D_1(T)$ and $D_2(T)$ are positive for $b_1\ge4$ and $T\ge0$, and all displayed coefficients in $P_1(T)$ and $P_2(T)$ are positive for $b_1\ge4$.  Hence \eqref{eq:corner-U-b} and \eqref{eq:corner-U-Vb} follow.

For the second range \eqref{eq:arithmetic-second-range}, Lemma \ref{lem:endpoint-principle} reduces the proof
to the four corners of $[1/c_1,V_{c_1}]\times[0,U]$.  At the two corners with
$v=0$, $\Kfun(u,0)=u\le V_{c_1}<L$.  The last inequality follows because, after writing $c_1=b_1(b_1-1)+T$, it becomes
$$
        \frac{(b_1+1)T^2+(2b_1^3-b_1)T+b_1^5-b_1^4-b_1^2-2b_1-1}{b_1(T+b_1^2-b_1+1)((b_1+1)T+b_1^3-b_1-1)}>0,
$$
which is positive for $b_1\ge2$ and $T\ge0$.
Thus, it remains to prove
\begin{equation}\label{eq:corner-Vc-U}
        \Kfun\left(\frac{1}{c_1},U\right)\le L,
        \qquad
        \Kfun(V_{c_1},U)\le L.
\end{equation}
For $b_1=2$, $c_1\geq3$, the two inequalities in \eqref{eq:corner-Vc-U} reduce respectively to
$$
        \frac{(c_1+2)(c_1^6+2c_1^5-16c_1^4+2c_1^3+43c_1^2-60c_1-36)}{2c_1(c_1-2)(c_1+1)(c_1+3)(c_1^3+2c_1^2-c_1+2)}>0
$$
and
$$
        \frac{9c_1^7+30c_1^6-134c_1^5-208c_1^4+621c_1^3+114c_1^2-1504c_1-272}{2(c_1-2)(c_1+1)(c_1+3)(3c_1-1)(3c_1^3+5c_1^2-4c_1+8)}>0,
$$
which holds for all $c_1\ge3$.
For $b_1=3$, $c_1\geq6$, the two inequalities in \eqref{eq:corner-Vc-U} reduce respectively to
$$
        \frac{c_1^7+4c_1^6-39c_1^5-68c_1^4+208c_1^3-90c_1^2-1080c_1-432}{3c_1(c_1-3)(c_1+1)(c_1+4)(c_1^3+2c_1^2-2c_1+6)}>0
$$
and
$$
        \frac{16c_1^7+56c_1^6-663c_1^5-806c_1^4+4031c_1^3-2265c_1^2-17226c_1-3807}{3(c_1-3)(c_1+1)(c_1+4)(4c_1-1)(4c_1^3+7c_1^2-9c_1+27)}>0,
$$
which holds for all $c_1\ge6$.
Finally let $b_1\ge4$, and put $c_1=b_1(b_1-1)+T$.  The two inequalities in \eqref{eq:corner-Vc-U} reduce respectively to
$$
        \frac{P_3(T)}{D_3(T)}>0,
        \qquad
        \frac{P_4(T)}{D_4(T)}>0,
$$
where
\begin{align*}
P_3(T)={}&T^7+(7b_1^2-7b_1+4)T^6+(21b_1^4-43b_1^3+44b_1^2-27b_1+6)T^5\\
        &+(35b_1^6-110b_1^5+165b_1^4-166b_1^3+103b_1^2-39b_1+4)T^4\\
        &+(b_1-1)(35b_1^7-115b_1^6+185b_1^5-219b_1^4+163b_1^3-90b_1^2\\
        &\qquad +24b_1-1)T^3\\
        &+b_1(21b_1^9-115b_1^8+290b_1^7-486b_1^6+598b_1^5-546b_1^4+380b_1^3\\
        &\qquad -188b_1^2+60b_1-6)T^2\\
        &+b_1^2(7b_1^{10}-47b_1^9+144b_1^8-289b_1^7+432b_1^6-494b_1^5\\
        &\qquad +445b_1^4-318b_1^3+170b_1^2-66b_1+12)T\\
        &+b_1^3(b_1^{11}-8b_1^{10}+29b_1^9-68b_1^8+119b_1^7-162b_1^6\\
        &\qquad +176b_1^5-158b_1^4+113b_1^3-66b_1^2+28b_1-8),\\[3pt]
D_3(T)={}&b_1(T+b_1^2+1)(T+b_1^2-2b_1)(T+b_1^2-b_1)(T+b_1^2-b_1+1)\\
        &\times\bigl(T^3+(3b_1^2-3b_1+2)T^2+(3b_1^4-6b_1^3+7b_1^2-5b_1+1)T\\
        &\qquad +b_1^6-3b_1^5+5b_1^4-6b_1^3+5b_1^2-2b_1\bigr),
\end{align*}
and
\begin{align*}
P_4(T)={}&(b_1+1)^2T^7+(b_1+1)(7b_1^3-3b_1+2)T^6\\
        &+b_1(21b_1^5-b_1^4-21b_1^3+6b_1^2-5b_1-11)T^5\\
        &+(35b_1^8-40b_1^7-20b_1^6+24b_1^5-37b_1^4+13b_1^2-7b_1-2)T^4\\
        &+(35b_1^{10}-80b_1^9+35b_1^8+6b_1^7-48b_1^6+55b_1^5+4b_1^4\\
        &\qquad -b_1^3+19b_1^2+5b_1-1)T^3\\
        &+b_1(21b_1^{11}-73b_1^{10}+81b_1^9-51b_1^8+8b_1^7+54b_1^6\\
        &\qquad -24b_1^5+19b_1^4+5b_1^3-11b_1^2+4)T^2\\
        &+b_1^2(7b_1^{12}-33b_1^{11}+57b_1^{10}-60b_1^9+49b_1^8-7b_1^7\\
        &\qquad -7b_1^6+5b_1^5-17b_1^4-11b_1^3-3b_1^2-7b_1-5)T\\
        &+b_1^3(b_1^{13}-6b_1^{12}+14b_1^{11}-20b_1^{10}+23b_1^9-17b_1^8\\
        &\qquad +9b_1^7-9b_1^6-2b_1^5-4b_1^4+b_1^3+5b_1+2),\\[3pt]
D_4(T)={}&b_1(T+b_1^2+1)(T+b_1^2-2b_1)(T+b_1^2-b_1+1)((b_1+1)T+b_1^3-b_1-1)\\
        &\times\bigl((b_1+1)T^3+(3b_1^3-b_1+1)T^2+(3b_1^5-3b_1^4+b_1^3-2b_1)T\\
        &\qquad +b_1^7-2b_1^6+2b_1^5-2b_1^4+b_1^3+b_1^2\bigr).
\end{align*}
All factors in $D_3(T)$ and $D_4(T)$ are positive for $b_1\ge4$ and $T\ge0$, and all displayed coefficients in $P_3(T)$ and $P_4(T)$ are positive for $b_1\ge4$.  Hence the two
inequalities in \eqref{eq:corner-Vc-U} hold.  This proves \eqref{eq:arithmetic-second-range} and completes the lemma.

\end{proof}

\end{document}
